\section{一般点への移行と有限降下}\label{rog:sec:generic}

\subsection{残余体stabilizationのEuler形式}

今度は$d=3$とし、$F\in k[x_1,\ldots,x_N]_3$を滑らかな射影cubicを定める三次式とする。
$K_F$を原点の残余体のstabilizationとする。
\begin{proposition}[非零のEulerベクトル]\label{rog:prop:euler}
$r=0,1,2$について
\[
 \chi_r:=\sum_p(-1)^p\dim_k\operatorname{Hom}(K_F,K_F(r)[p])
 =\sum_{\substack{0\le s\le N\\s\equiv-r\pmod3}}(-1)^s\binom Ns.
 \tag{RO3}\label{rog:eq:euler}
\]
この三成分ベクトルは零でない。より正確に
\[
 \chi_0+\chi_1+\chi_2=0,\qquad
 \chi_0^2+\chi_1^2+\chi_2^2=2\cdot3^{N-1}.
 \tag{RO4}\label{rog:eq:norm}
\]
$N=5$では$(\chi_0,\chi_1,\chi_2)=(-9,9,0)$である。
\end{proposition}
\begin{proof}
$F=\sum_i x_if_i$、$f_i=(\partial_iF)/3$、
$f_i=\sum_jx_ja_{ij}$、$a_{ij}=(\partial_jf_i)/2$と置く。
外積作用$e_i\wedge-$の二重次数を$(-1,1)$、
収縮$\iota_i$を$(1,-1)$、周期変数$t$を$(2,-3)$とする。
$K_F$の微分は
\[
 Q=\sum_i x_i\iota_i+t\sum_i f_i(e_i\wedge-),\qquad Q^2=tF.
\]
そのendomorphism代数のsuper交換子微分について
\[
 d(e_i\wedge-)=x_i,\qquad
 \xi_i:=\iota_i-t\sum_j a_{ij}(e_j\wedge-),\qquad d\xi_i=0.
\]
順序を固定した$e_J\xi_I$はendomorphism加群の基底である。
実際、標準の$e_J\iota_I$からの変換は三角的である。
この基底では微分は$x_1,\ldots,x_N$のKoszul微分となる。
また$\xi_i\xi_j+\xi_j\xi_i=-t(a_{ij}+a_{ji})$であり、右辺は原点で零となる一次式なので境界である。
従って、周期変数を残した二重次数付きコホモロジーは
\[
 \bigwedge(\xi_1,\ldots,\xi_N)\otimes_k k[t,t^{-1}],
 \qquad\deg\xi_i=(1,-1)
 \tag{RO5}\label{rog:eq:exterior}
\]
となる。これは\cite[\S5.6]{ROGDyckerhoff}のClifford計算の三次の場合でもある。

$\xi_I t^j$が内部shift $r$のHomに寄与する条件は
$-|I|-3j=r$であり、cohomological次数は$|I|+2j$である。
$0\le|I|\le N$なので各$r$について有限個しか寄与せず、符号は$(-1)^{|I|}$。
これで\eqref{rog:eq:euler}を得る。

$A(z)=(1-z)^N$を$z^3-1$で割った余りの係数が、順序を除いて$\chi_r$である。
原始三乗根$\zeta$において$A(\zeta)=(1-\zeta)^N\ne0$なので三成分は同時には零でない。
$A(1)=0$、$|1-\zeta|^2=3$と三点のFourier公式から\eqref{rog:eq:norm}を得る。
これは整数恒等式なので、基礎体に依らない。$N=5$の値は直接代入する。
\end{proof}

\begin{remark}
一般$N$で自己Euler数$\chi_0$だけを使うのは誤りである。
例えば$N=9$では$\chi_0=0$となるが、三成分ベクトルは非零である。
また上の群は\emph{生成対象間のHom}であり、$\operatorname{Nat}(\mathrm{Id},\rho^r)$ではない。
\end{remark}

\subsection{固定した有限対象上では形式functorを先に局所化できる}

$R=\mathbb C[h]_{(h)}$、$K=\mathbb C(h)$、
$\widehat R=\mathbb C[[h]]$、$L=\mathbb C((h))$とする。
代数的中心三次対$(A_R,\Omega_R)$について、次の既存入力を使用する。
\begin{description}
\item[\textup{(L)}] 次数別flatな局所量子化と有限Koszul分解、
線束生成残余圏のproper模型、regular底のstabilization/window比較がある。
原点の$R$-加群の有限rank stabilizationを$K_R$とすると
\[
 T_R=K_R\oplus K_R(1)\oplus K_R(2),\quad
 \Gamma_R=R\operatorname{End}(T_R),\quad
 \mathcal K^{\rm lin}_{B_R}\simeq\operatorname{Perf}(\Gamma_R),
 \qquad \Gamma_R\in\operatorname{Perf}(R).
 \tag{RO6}\label{rog:eq:source}
\]
この模型とHomは$\widehat R,L$への平坦底変更に両立する。
\item[\textup{(CF)}] 次数別完備底上の有限rank graded MFに、
内部shiftと指定周期性を保つ実際の$\widehat R$-線形$A_\infty$同値
\[
 \Phi:\operatorname{MF}^{\mathbb Z}(A^{\rm gr}_{\widehat R},\Omega)
 \simeq\operatorname{MF}^{\mathbb Z}(\widehat R[x],W_h)
 \tag{RO7}\label{rog:eq:completed}
\]
がある。$W_h\bmod h=F$は滑らかで、$\Phi$は剰余で恒等。
単なる両立したArtin同値の列ではなく、
定理\ref{cft:thm:complete}で構成した固定内部次数内で収束する
実際の関手を用いる。
\end{description}
記号(CF)は定理\ref{cft:thm:complete}で構成した完備底上の関手を指す。
本節では、これらの比較定理を入力として、
固定三対象の局所化から一般点の生成性へ進む部分を独立に記述する。



\begin{lemma}[有限対象への制限とLaurent局所化]\label{rog:lem:localization}
\textup{(L),(CF)}の下で$E=\Phi(K_{\widehat R})$と置くと、
$E$は通常の$\widehat R[x]$上の有限rank graded MFであり、
\[
 \operatorname{Perf}(\Gamma_R\otimes_RL)
 \simeq\operatorname{thick}\{E_L,E_L(1),E_L(2)\}
 \tag{RO8}\label{rog:eq:image}
\]
というrotationと周期性を保つ同値がある。
また$\operatorname{Hom}(E,E(r))$は$r=0,1,2$で$\widehat R$-perfectであり、
その特殊・一般ファイバー間のEuler数は\eqref{rog:eq:euler}に等しい。
\end{lemma}
\begin{proof}
固定した有限rankの基礎加群と固定cohomological次数では、MF行列の各成分の
内部次数と$t$の指数は固定される。従って各成分は有限自由な
$\widehat R$-加群$\widehat R[x]_j$に属する。
(CF)の$F_{2q+1}=O(h^q)$によるtwistingの和は、その有限自由成分内で収束する。
結果は$x$の次数を無制限に増やす級数ではなく、$h$の形式係数を持つ有限次数の多項式行列である。
固定個数の射に対する各Taylor成分も同じ理由で通常の$\widehat R$-多重線形写像となる。
恒等式と逆functorは全Artin剰余で成立し、成分の分離性によって$\widehat R$上で成立する。

従って$T_{\widehat R}$の有限対象部分圏では、(CF)は\emph{通常の}endomorphism複体の
$A_\infty$準同型を与える。$\widehat R\to L$は平坦なので、これをテンソルすると
通常の準同型である。固定cohomological次数のHomの形成はこの平坦底変更と可換する。
$\Phi$を任意のLaurent係数対象へ新たに拡張することは使わず、
この固定部分圏のpretriangulated・冪等包絡を取って\eqref{rog:eq:image}を得る。
有限twisted complexのtwisting行列は厳密上三角に取れるため、包絡への延長では有限個のTaylor成分だけが使われる。
したがって$L$への移行後に新しい無限級数の収束を要求しない。

Homのperfect性は(L)と(CF)から従い、$K_R(r)$は周期性により$T_R$のshiftの直和因子である。
有界有限射影複体のEuler数は特殊化で不変なので、剰余の$K_F$について計算した
\eqref{rog:eq:euler}が一般ファイバーでも保たれる。
内部shiftと周期性の自然比較も、この固定部分圏への制限と平坦底変更で保たれる。
\end{proof}

\subsection{体拡大上の同値から有限代数拡大上の同値へ}

この降下の基礎は、smooth proper dg代数の有限表示・降下
\cite[Th\'eor\`eme~1]{ROGToenDefinition}、perfect加群のfiltered colimitとの両立
\cite[Lemma~2.10]{ROGTV}、および同値の核を単位・余単位によって判定する原理
\cite[Lemma~3.25]{ROGTV}である。
以下では片側のsmoothnessから始め、
rotationと指定した有限個の周期図式を保つ形で、これらの原理を適用する。
そのため、降下核と同値を判定するconeの底上のperfect性を明示する。

\begin{lemma}[固定された可換模型へのgeneric降下]\label{rog:lem:descent}
$K$を標数零の体、$L/K$を任意の体拡大、$\Gamma,\mathsf C$を$K$上properなdg代数とし、
$\mathsf C$をsmoothとする。両側にrotationを表す可逆双加群と、必要なら三次周期同型を指定する。
$L$上でrotation・周期性を保つMorita同値があれば、ある有限代数拡大$K'/K$上でも
同じ二つの代数の間にそれらを保つMorita同値がある。
\end{lemma}
\begin{proof}
$L$上の同値を可換側から非可換側への核$P_L$で表す。
これは右$\Gamma_L$-perfectであり、$\mathsf C_L$のsmoothnessにより
$\mathsf C_L^{\rm op}\otimes_L\Gamma_L$上でもperfectである
\cite[Lemma~2.8]{ROGTV}。
$L$を有限生成$K$-部分代数のfiltered colimitとして表し、
\cite[Lemma~2.10]{ROGTV}により核とrotation比較を一つの有限生成整域$D\subset L$へ降ろす。
$\mathsf C_D$は$D$-properなので、包絡代数上のperfect加群を右$\Gamma_D$へ制限してもperfectである。

$Q_D=R\operatorname{Hom}_{\Gamma_D}(P_D,\Gamma_D)$と置く。
$P_D$が右$\Gamma_D$-perfectで、$\Gamma_D$が$D$-properなので、
$P_D$と$Q_D$はいずれも$D$-perfectである。
$Q_D$を右$\mathsf C_D$-加群と見ると、$\mathsf C_D$のsmoothnessと
\cite[Lemma~2.8(2)]{ROGTV}により右$\mathsf C_D$-perfectでもある。
従って$-\otimes_{\mathsf C_D}^{\mathbf L}P_D$と
$-\otimes_{\Gamma_D}^{\mathbf L}Q_D$はperfect加群の圏の間の随伴対をなし、
その単位・余単位は双加群の射
\[
 \mathsf C_D\longrightarrow P_D\otimes_{\Gamma_D}^{\mathbf L}Q_D,
 \qquad Q_D\otimes_{\mathsf C_D}^{\mathbf L}P_D\longrightarrow\Gamma_D
\]
で表される。第一のテンソル積は$P_D$の右$\Gamma_D$-perfect性と
$Q_D$の$D$-perfect性から、第二のテンソル積は$Q_D$の右$\mathsf C_D$-perfect性と
$P_D$の$D$-perfect性から、それぞれ$D$-perfectである。
両方の対角も$D$-perfectなので、単位・余単位のconeは$D$-perfectとなる。
また$P_D$の右$\Gamma_D$-perfect性により、$Q_D$とこの随伴の形成は
$D\to L$の底変更と可換する。従って両coneは$L$への底変更後に零となる。

$D$はnoetherian整域なので、これらのconeは有限個の有限生成cohomology加群を持つ。
$D\subset L$より各加群は$D$-torsionであり、有限生成元のannihilatorの積を取れば、
ある$0\ne f\in D$の反転後に両coneが零となる。
rotation比較の両辺も、可逆双加群によるテンソルが片側のperfect性を保つことと
両代数のproper性から$D$-perfectであり、そのconeを同じ方法で消せる。
指定周期性の図式はperfect双加群間の有限個の射の等式であるから、
必要なら段階を拡大し、さらに一つの元を反転して等式を保つ。
任意の高次群作用を有限の周期図式と同一視してはいない。

$\operatorname{Spec}D_f$の閉点を一つ取る。その剰余体$K'$はZariskiの補題により$K$の有限拡大であり、
標数零なので分離的である。そこでの底変更が所要の同値を与える。
これは元の$L$上の形式分枝を全次数で再現する主張ではない。
\end{proof}

\begin{theorem}[一変数量子化の幾何学的一般点]\label{rog:thm:geometric-generic}
\textup{(L),(CF)}を満たす一変数中心三次量子化に対し、
ある有限拡大$K'/\mathbb C(h)$と滑らかな三次超曲面$Y/K'$が存在して
\[
 (\mathcal K^{\rm lin}_{B_K}\otimes_KK',\rho)
 \simeq(\operatorname{Ku}(Y),\rho_Y)
 \tag{RO10}\label{rog:eq:generic}
\]
となる。$\mathcal K^{\rm lin}$と全coherent残余圏との同一性は、この定理に含めない。
\end{theorem}
\begin{proof}
補題\ref{rog:lem:localization}により、$L=\mathbb C((h))$上で像は
$E_L,E_L(1),E_L(2)$の厚い閉包になる。
命題\ref{rog:prop:euler}とproper HomのEuler数の底変更により$E_L\ne0$である。
$W_h$の判別式は特殊ファイバーで非零なので$L$上も滑らかである。
$L$の代数閉包$\overline L$へ移し、命題\ref{gcd:prop:bound}を適用すると
像は可換MF圏全体となる。体上のOrlov比較により$\overline L$上で可換
Kuznetsov成分との同値を得る。以下では$L$をこの$\overline L$に置き換えて
有限降下を行う。
$W_h$の有限個の三次係数と判別式の逆を含む有限生成$K$-部分代数$D_0\subset L$を取ると、
$Y_0=V(W_{D_0})$は$D_0$上の滑らかな射影cubic族である。
その線束の有限和を残余成分へ射影した生成子を$H_{D_0}$、
生成dg代数を$\mathsf C_{D_0}=R\operatorname{End}(H_{D_0})$とする。
これはsmooth properで、底変更と両立する。
ここで生成子は相対ampleな全線束が生成することと、座標のKoszul完全列によって
$N$個の連続したtwistへ減らすことから得られる。
射影は有限個のexceptional mutationによる。
この有限線束生成の基底環に関する一般性は\cite[\S36.16]{ROGStacks}を参照する。

$L$上の同値核を、$D_0$を含む有限生成$K$-部分代数へ降ろす。
以後は補題\ref{rog:lem:descent}の証明と同じであり、
properな二つのconeを消して閉点を選ぶ。
その点では$Y_0$も滑らかなcubicへ特殊化し、\eqref{rog:eq:generic}を得る。
この一般存在定理は、$Y$の係数の閉公式や同値核の行列を明示するものではない。
\end{proof}

